\chapter{Vaginal problems}
\index{Vaginal problems}

\section*{Definition and Overview}
Vaginal problems are a common reason for medical consultation and cover a broad range of symptoms, including abnormal discharge, itching, soreness, dryness, unusual bleeding, and pelvic discomfort. The term is deliberately wide because the underlying causes are many: infections, hormonal changes, allergic reactions, skin conditions, structural changes, and, less commonly, growths or cancer. Because symptoms overlap considerably between benign and more serious conditions, an accurate history and examination are the cornerstones of safe management.

The healthy vagina maintains its own defence system. Oestrogen keeps the lining moist and thick, while \textit{Lactobacillus} bacteria produce acid that suppresses harmful organisms. Many "problems" are simply variations of normal physiology, such as the increased clear, stretchy discharge around ovulation or the heavier, thicker discharge before a period. Distinguishing a normal variant from a condition requiring treatment is the first step in managing any vaginal symptom. This chapter takes a symptom-based approach and points to the more detailed chapters on specific infections elsewhere in this book.

\section*{Epidemiology}
Vaginal symptoms are among the most common reasons women visit a general practitioner, and most women experience at least one significant vaginal symptom during their lifetime. Abnormal discharge and vulval itching are the most frequent complaints, and their causes vary with age. During the reproductive years, infection is the most common cause, with bacterial vaginosis, thrush, and sexually transmitted infections (STIs) heading the list. Around one-third of women report a vulvovaginal complaint in any given year.

After the menopause, hormonal change becomes the dominant theme: oestrogen levels fall, the vaginal lining thins and becomes drier, and the pH rises. Atrophic vaginitis affects a large proportion of post-menopausal women, often beginning during the perimenopause. Structural problems such as pelvic organ prolapse also become more common with age and after childbirth. In girls before puberty, the vagina is thin and less protective, so vulvovaginitis from irritants or bowel organisms is a frequent presentation, whereas the hormonal infections of adult women are uncommon.

\section*{Aetiology and Pathophysiology}
Vaginal problems have several broad groups of causes, and often more than one factor is at work.

\begin{itemize}
    \item \textbf{Infections:} bacterial vaginosis, vulvovaginal candidiasis (thrush), trichomoniasis, and STIs such as chlamydia and gonorrhoea; in girls, non-specific vulvovaginitis related to poor hygiene or irritants is typical
    \item \textbf{Hormonal changes:} low oestrogen around the menopause or after childbirth produces dryness, thinning, and soreness; hormonal contraceptives and pregnancy alter discharge and infection risk
    \item \textbf{Irritants and allergies:} soaps, bubble bath, perfumed products, panty liners, laundry detergents, spermicides, and synthetic underwear can cause contact dermatitis of the vulva
    \item \textbf{Skin conditions:} eczema, psoriasis, and lichen sclerosus commonly affect the vulval skin and cause itching and soreness
    \item \textbf{Structural change:} pelvic organ prolapse produces a feeling of a lump, heaviness, or dragging; a urethral caruncle or cyst can cause a visible swelling
    \item \textbf{Foreign bodies:} retained tampons, condoms, or sex toys cause discharge, odour, and occasionally serious infection
    \item \textbf{Growths:} cervical polyps are common and benign, whereas persistent abnormal bleeding or discharge can occasionally signal pre-cancer or cancer of the cervix, vagina, or vulva
    \item \textbf{Pain disorders:} vulvodynia and vestibulodynia cause chronic burning or pain without any visible abnormality
\end{itemize}

Understanding normal discharge is the first step in recognising a problem. Around ovulation, oestrogen peaks and discharge becomes clear, stretchy, and mucus-like; in the second half of the cycle it is often thicker and whiter; and in pregnancy it generally increases. A discharge only needs investigation when it changes character or smell, causes itching, soreness, or burning, is blood-stained, or is accompanied by pelvic pain or fever. Likewise, light spotting after sex or at ovulation can be normal, but repeated or heavy bleeding outside a period is not, and bleeding after the menopause always requires review.

\section*{Clinical Features}
\begin{itemize}
    \item \textbf{Discharge:} a change in amount, colour, consistency, or smell of vaginal discharge; note whether it is itchy, frothy, curdy, fishy, or blood-stained
    \item \textbf{Itching and soreness:} irritation of the vulva and vagina, which may be worse at night, after sex, or with washing
    \item \textbf{Dryness and discomfort:} common around the menopause, causing soreness with sex and reduced lubrication
    \item \textbf{Bleeding:} bleeding between periods, after sex, or after the menopause is never normal and requires investigation
    \item \textbf{Lump or bulging:} a feeling of a lump in the vagina, pelvic heaviness, or something "coming down" suggests prolapse
    \item \textbf{Pain:} pain with sex, during urination, or deep pelvic aching may accompany infection, skin disease, or pelvic disorders
    \item \textbf{Red flags:} fevers, severe pelvic pain, heavy bleeding, or symptoms that persist despite treatment should always be taken seriously
\end{itemize}

\section*{Differential Diagnosis}
Because vaginal symptoms overlap, the differential diagnosis is wide.

\begin{itemize}
    \item Normal physiological discharge and its cyclical variations
    \item Bacterial vaginosis and vaginal thrush
    \item Sexually transmitted infections, including trichomoniasis, chlamydia, and gonorrhoea
    \item Atrophic vaginitis and vulvovaginal dryness related to the menopause
    \item Contact dermatitis and allergic reactions to intimate products
    \item Vulval skin diseases, including eczema, psoriasis, and lichen sclerosus
    \item Pelvic organ prolapse (cystocele, rectocele, or uterine prolapse)
    \item Vulvodynia and other chronic pain syndromes
    \item Cervical polyps, endometriosis, and pelvic inflammatory disease
    \item Malignancy of the cervix, vulva, or vagina, particularly with persistent abnormal bleeding, a lump, or unexplained symptoms
\end{itemize}

\section*{When to Seek Medical Care}
Any new, persistent, or worrying vaginal symptom should be discussed with a doctor, nurse, or sexual health clinic. Seek review if a discharge is offensive, blood-stained, or associated with pelvic pain, if itching or soreness interferes with daily life, if there is bleeding between periods or after sex, or if symptoms return repeatedly despite treatment. A lump, a feeling of prolapse, or symptoms after the menopause warrant particular attention.

\textbf{Contact your local emergency services for:}
\begin{itemize}
    \item Severe or one-sided lower abdominal pain, especially with fever
    \item Heavy vaginal bleeding that soaks pads rapidly
    \item Fainting, dizziness, or signs of shock
    \item Sudden high fever with a sunburn-like rash and vomiting, which may indicate toxic shock syndrome from a retained tampon
\end{itemize}

\section*{Diagnosis and Investigations}
\begin{itemize}
    \item \textbf{Detailed history:} character and timing of symptoms, menstrual and pregnancy history, sexual history, contraceptives, medicines, and self-care tried
    \item \textbf{Vaginal and cervical examination:} inspection of the vulva, vagina, and cervix with speculum, always performed with consent and explanation
    \item \textbf{Swabs and tests:} pH testing, microscopy, and molecular (PCR) testing for thrush, bacterial vaginosis, trichomoniasis, chlamydia, and gonorrhoea
    \item \textbf{Cervical screening:} regular cervical screening tests remain the best protection against cervical pre-cancer and cancer
    \item \textbf{Imaging:} pelvic ultrasound may be arranged for pain, bleeding, or suspected pelvic masses
    \item \textbf{Biopsy:} a small skin sample may be needed to confirm lichen sclerosus or other skin conditions
    \item \textbf{Specialist referral:} gynaecology review for persistent bleeding, suspected prolapse, chronic pain, or findings that do not respond to first-line care
\end{itemize}

\section*{Management}
Management is directed at the underlying cause.

\begin{itemize}
    \item \textbf{Infections:} treated with the appropriate antibiotic or antifungal, as outlined in the chapters on vaginal infection and vaginal thrush
    \item \textbf{Dryness and atrophic vaginitis:} vaginal moisturisers and lubricants, and topical or low-dose vaginal oestrogen for post-menopausal women, which is safe for most and highly effective
    \item \textbf{Irritation and allergy:} stop using perfumed soaps, bubble bath, and other products; use plain water or a mild unperfumed wash and cotton underwear
    \item \textbf{Skin conditions:} steroid creams or ointments under medical supervision, with long-term follow-up for lichen sclerosus because of its small cancer risk
    \item \textbf{Prolapse:} pelvic floor exercises, lifestyle measures to avoid straining, a pessary, or surgery depending on severity and the woman's preferences
    \item \textbf{Polyps and growths:} simple removal of polyps; persistent or suspicious lesions are investigated further
    \item \textbf{Pain disorders:} a multidisciplinary approach including topical treatments, physiotherapy, counselling, and pain management is most effective for vulvodynia
    \item \textbf{Symptom control:} cool compresses, plain emollients, and avoiding tight clothing can relieve discomfort while a diagnosis is being made
\end{itemize}

Many women worry that their symptoms are embarrassing or trivial and delay seeking help. It may help to know that vaginal symptoms are routine presentations in general practice, are discussed non-judgementally every day, and are almost never a sign of poor hygiene or sexual behaviour. A clinician's first questions will cover the timing and nature of symptoms, the menstrual cycle, contraception, and any sexual activity, simply because these are the details that point to the cause. Treatment often begins the same day, and most women leave a consultation with both a diagnosis and a plan.

\section*{Complications}
\begin{itemize}
    \item Untreated infections can spread to the pelvis, causing pelvic inflammatory disease, chronic pain, and reduced fertility
    \item Recurrent or chronic symptoms can significantly impair quality of life, sleep, and sexual relationships
    \item Delayed diagnosis of lichen sclerosus allows scarring (loss of vulval architecture) and slightly increases cancer risk
    \item Untreated pelvic organ prolapse can worsen with time, causing urinary leakage, difficulty emptying, and ulceration
    \item Persistent post-menopausal bleeding carries a risk of gynaecological malignancy and must never be dismissed
    \item Vulvodynia can become entrenched and disabling without early, sympathetic, multidisciplinary care
\end{itemize}

\section*{Prognosis and Outcomes}
Most vaginal problems resolve completely with correct treatment. Infections clear within days, and simple measures such as stopping irritants or beginning moisturisers often produce rapid improvement. Atrophic vaginitis responds well to vaginal oestrogen and can be managed effectively for years. Even prolapse, chronic skin disease, and pain disorders can be substantially improved with a combination of treatments and specialist support.

The main risk to long-term health is delay: infections that ascend to the pelvis, and pre-cancerous or cancerous changes detected late, carry far worse outcomes. This is why persistent symptoms, unusual bleeding, and post-menopausal change always deserve prompt assessment. With early review and appropriate treatment, the outlook for nearly every vaginal problem is excellent.

It is also worth emphasising that a symptom does not have to be severe to deserve attention. Mild but persistent itching, a discharge that will not settle, or discomfort that has been accepted for months are all worth a discussion. Many women endure symptoms for long periods because they assume they will resolve or feel they should not bother their doctor, yet most are easily treated once accurately diagnosed. The simplest advice is that any symptom which changes, persists, or interferes with daily life or relationships is a reason to seek review, and that treatment is available and effective at every age.

\section*{Prevention and Self-Care}
\begin{itemize}
    \item Avoid douching, vaginal deodorants, and perfumed intimate products
    \item Wash the vulva with plain water or a mild unperfumed cleanser and pat dry gently
    \item Wear cotton underwear, avoid tight synthetic clothing, and change wet swimwear promptly
    \item Use condoms to reduce the risk of sexually transmitted infections
    \item Use a lubricant if dryness causes discomfort with sex
    \item Maintain regular cervical screening appointments
    \item Practise pelvic floor exercises, especially after childbirth and as menopause approaches
    \item Keep a symptom diary if symptoms are recurrent, noting their timing around the cycle, to help identify triggers
\end{itemize}

\section*{Sources and Further Reading}
The following organisations provide reliable information on vaginal problems and their management.

\begin{itemize}
    \item NHS. \textit{Information on vaginal discharge, dryness, and related symptoms.} https://www.nhs.uk/conditions/
    \item Mayo Clinic. \textit{Guide to vaginitis and women's health symptoms.} https://www.mayoclinic.org/
    \item MedlinePlus. \textit{Health information on vaginal diseases and women's health.} https://medlineplus.gov/
    \item MSD Manual. \textit{Patient education on vaginal symptoms and their causes.} https://www.msdmanuals.com/
    \item Jean Hailes for Women's Health. \textit{Resources on vulval and vaginal health.} https://www.jeanhailes.org.au/
    \item HealthDirect Australia. \textit{Consumer information on vaginal symptoms and women's health.} https://www.healthdirect.gov.au/
\end{itemize}
